% ===================================================================
%  Общая преамбула для семинаров по эконометрике.
%  Подключается из <seminar_NN>_problems.tex и <seminar_NN>_solutions.tex
%  строкой % ===================================================================
%  Общая преамбула для семинаров по эконометрике.
%  Подключается из <seminar_NN>_problems.tex и <seminar_NN>_solutions.tex
%  строкой % ===================================================================
%  Общая преамбула для семинаров по эконометрике.
%  Подключается из <seminar_NN>_problems.tex и <seminar_NN>_solutions.tex
%  строкой % ===================================================================
%  Общая преамбула для семинаров по эконометрике.
%  Подключается из <seminar_NN>_problems.tex и <seminar_NN>_solutions.tex
%  строкой \input{../_tex/econ_preamble}
%
%  Компилировать XeLaTeX'ом (нужен fontspec), два прохода.
%  Драйвер обязан определить \newif\ifsolutions ДО этого \input.
% ===================================================================

\usepackage{fontspec}
\setmainfont{Times New Roman}
\setsansfont{Arial}
\setmonofont{Consolas}

\usepackage[russian]{babel}
\usepackage{amsmath,amssymb}
\usepackage{booktabs,array,multirow}
\usepackage{enumitem}
% Регистрируем русский счётчик, иначе метка label=\asbuk*) не работает.
\makeatletter
\AddEnumerateCounter{\asbuk}{\@asbuk}{щ}
\makeatother
\usepackage{xcolor}
\usepackage{framed}
\usepackage{comment}
\usepackage[a4paper,margin=25mm,headheight=15pt]{geometry}
\usepackage{fancyhdr}
\usepackage[hidelinks]{hyperref}
\usepackage{titlesec}
\usepackage{tikz}
\usetikzlibrary{calc}

% ------------------------- цвета -----------------------------------
\definecolor{accent}{RGB}{0,80,130}
\definecolor{solbar}{RGB}{0,120,90}
\definecolor{notebar}{RGB}{170,120,0}
\definecolor{xlbar}{RGB}{110,60,140}

% --------------- название семинара для колонтитула ------------------
\newcommand{\seminarname}{Семинар}
\newcommand{\setseminar}[1]{\renewcommand{\seminarname}{#1}}

% ------------------------ колонтитулы -------------------------------
\pagestyle{fancy}
\fancyhf{}
\fancyhead[L]{\small\itshape Эконометрика \textperiodcentered\ \seminarname}
\fancyhead[R]{\small\itshape \ifsolutions Решения\else Условия\fi}
\fancyfoot[C]{\small\thepage}
\renewcommand{\headrulewidth}{0.4pt}

% ------------------------ заголовки ---------------------------------
\titleformat{\section}{\large\bfseries\color{accent}}{\thesection.}{0.6em}{}
\titleformat{\subsection}{\bfseries}{\thesubsection.}{0.6em}{}
\titlespacing*{\section}{0pt}{1.4\baselineskip}{0.6\baselineskip}

% --------------------- шапка документа ------------------------------
% \seminartitle{Семинар 2}{подзаголовок}
\newcommand{\seminartitle}[2]{%
  \setseminar{#1}%
  \begin{center}
    {\LARGE\bfseries\color{accent} #1}\\[2pt]
    {\large #2}\\[4pt]
    {\itshape \ifsolutions Условия и решения \else Условия задач \fi}
  \end{center}
  \vspace{0.8\baselineskip}}

% --------------------- блок «Задача N» ------------------------------
\newcommand{\problem}[1]{%
  \bigskip\par\noindent
  {\bfseries\color{accent}\large Задача #1}\par\nopagebreak\smallskip}

% ------------------- вертикальная линейка слева ----------------------
\newenvironment{leftbarred}[1][solbar]{%
  \def\FrameCommand{{\color{#1}\vrule width 2.5pt}\hspace{10pt}}%
  \MakeFramed{\advance\hsize-\width \FrameRestore}}%
 {\endMakeFramed}

% ------------------------ «Замечание» --------------------------------
\newenvironment{remark}{%
  \par\smallskip\begin{leftbarred}[notebar]%
  \noindent{\bfseries\color{notebar}Замечание.}\ \itshape}%
 {\end{leftbarred}\par\smallskip}

% ------------------------- «В Excel» ---------------------------------
% Видно и в условиях, и в решениях: это часть задания, а не ответ.
\newenvironment{excel}{%
  \par\smallskip\begin{leftbarred}[xlbar]%
  \noindent{\bfseries\color{xlbar}В Excel.}\ }%
 {\end{leftbarred}\par\smallskip}

% Имя пункта меню / кнопки в интерфейсе Excel
\newcommand{\ui}[1]{\textsf{\small #1}}
% Формула Excel
\newcommand{\xl}[1]{\texttt{\small #1}}

% ------------------------- «Решение» ---------------------------------
% Если \solutionsfalse — всё содержимое окружения просто выбрасывается.
\ifsolutions
  \newenvironment{solution}{%
    \par\smallskip\begin{leftbarred}[solbar]%
    \noindent{\bfseries\color{solbar}Решение.}\ }%
   {\end{leftbarred}\par\smallskip}
\else
  \excludecomment{solution}
\fi

% ---------------------- сокращения для формул ------------------------
\newcommand{\TSS}{\mathrm{TSS}}
\newcommand{\ESS}{\mathrm{ESS}}
\newcommand{\RSS}{\mathrm{RSS}}
% Матожидание и дисперсия — ажурным шрифтом (blackboard bold, amssymb).
\newcommand{\E}{\mathbb{E}}   % математическое ожидание
\newcommand{\D}{\mathbb{D}}   % дисперсия
\newcommand{\V}{\D}           % то же; старое имя, им набраны тексты семинаров
\newcommand{\Cov}{\mathrm{Cov}}
\newcommand{\yh}{\widehat{Y}}
\newcommand{\bh}{\widehat{\beta}}
\newcommand{\ah}{\widehat{\alpha}}
\renewcommand{\arraystretch}{1.25}

% Небольшой запас на растяжку строк: убирает единичные overfull-строки
% в русском тексте, где переносов меньше, чем в английском.
\emergencystretch=1.5em

%
%  Компилировать XeLaTeX'ом (нужен fontspec), два прохода.
%  Драйвер обязан определить \newif\ifsolutions ДО этого \input.
% ===================================================================

\usepackage{fontspec}
\setmainfont{Times New Roman}
\setsansfont{Arial}
\setmonofont{Consolas}

\usepackage[russian]{babel}
\usepackage{amsmath,amssymb}
\usepackage{booktabs,array,multirow}
\usepackage{enumitem}
% Регистрируем русский счётчик, иначе метка label=\asbuk*) не работает.
\makeatletter
\AddEnumerateCounter{\asbuk}{\@asbuk}{щ}
\makeatother
\usepackage{xcolor}
\usepackage{framed}
\usepackage{comment}
\usepackage[a4paper,margin=25mm,headheight=15pt]{geometry}
\usepackage{fancyhdr}
\usepackage[hidelinks]{hyperref}
\usepackage{titlesec}
\usepackage{tikz}
\usetikzlibrary{calc}

% ------------------------- цвета -----------------------------------
\definecolor{accent}{RGB}{0,80,130}
\definecolor{solbar}{RGB}{0,120,90}
\definecolor{notebar}{RGB}{170,120,0}
\definecolor{xlbar}{RGB}{110,60,140}

% --------------- название семинара для колонтитула ------------------
\newcommand{\seminarname}{Семинар}
\newcommand{\setseminar}[1]{\renewcommand{\seminarname}{#1}}

% ------------------------ колонтитулы -------------------------------
\pagestyle{fancy}
\fancyhf{}
\fancyhead[L]{\small\itshape Эконометрика \textperiodcentered\ \seminarname}
\fancyhead[R]{\small\itshape \ifsolutions Решения\else Условия\fi}
\fancyfoot[C]{\small\thepage}
\renewcommand{\headrulewidth}{0.4pt}

% ------------------------ заголовки ---------------------------------
\titleformat{\section}{\large\bfseries\color{accent}}{\thesection.}{0.6em}{}
\titleformat{\subsection}{\bfseries}{\thesubsection.}{0.6em}{}
\titlespacing*{\section}{0pt}{1.4\baselineskip}{0.6\baselineskip}

% --------------------- шапка документа ------------------------------
% \seminartitle{Семинар 2}{подзаголовок}
\newcommand{\seminartitle}[2]{%
  \setseminar{#1}%
  \begin{center}
    {\LARGE\bfseries\color{accent} #1}\\[2pt]
    {\large #2}\\[4pt]
    {\itshape \ifsolutions Условия и решения \else Условия задач \fi}
  \end{center}
  \vspace{0.8\baselineskip}}

% --------------------- блок «Задача N» ------------------------------
\newcommand{\problem}[1]{%
  \bigskip\par\noindent
  {\bfseries\color{accent}\large Задача #1}\par\nopagebreak\smallskip}

% ------------------- вертикальная линейка слева ----------------------
\newenvironment{leftbarred}[1][solbar]{%
  \def\FrameCommand{{\color{#1}\vrule width 2.5pt}\hspace{10pt}}%
  \MakeFramed{\advance\hsize-\width \FrameRestore}}%
 {\endMakeFramed}

% ------------------------ «Замечание» --------------------------------
\newenvironment{remark}{%
  \par\smallskip\begin{leftbarred}[notebar]%
  \noindent{\bfseries\color{notebar}Замечание.}\ \itshape}%
 {\end{leftbarred}\par\smallskip}

% ------------------------- «В Excel» ---------------------------------
% Видно и в условиях, и в решениях: это часть задания, а не ответ.
\newenvironment{excel}{%
  \par\smallskip\begin{leftbarred}[xlbar]%
  \noindent{\bfseries\color{xlbar}В Excel.}\ }%
 {\end{leftbarred}\par\smallskip}

% Имя пункта меню / кнопки в интерфейсе Excel
\newcommand{\ui}[1]{\textsf{\small #1}}
% Формула Excel
\newcommand{\xl}[1]{\texttt{\small #1}}

% ------------------------- «Решение» ---------------------------------
% Если \solutionsfalse — всё содержимое окружения просто выбрасывается.
\ifsolutions
  \newenvironment{solution}{%
    \par\smallskip\begin{leftbarred}[solbar]%
    \noindent{\bfseries\color{solbar}Решение.}\ }%
   {\end{leftbarred}\par\smallskip}
\else
  \excludecomment{solution}
\fi

% ---------------------- сокращения для формул ------------------------
\newcommand{\TSS}{\mathrm{TSS}}
\newcommand{\ESS}{\mathrm{ESS}}
\newcommand{\RSS}{\mathrm{RSS}}
% Матожидание и дисперсия — ажурным шрифтом (blackboard bold, amssymb).
\newcommand{\E}{\mathbb{E}}   % математическое ожидание
\newcommand{\D}{\mathbb{D}}   % дисперсия
\newcommand{\V}{\D}           % то же; старое имя, им набраны тексты семинаров
\newcommand{\Cov}{\mathrm{Cov}}
\newcommand{\yh}{\widehat{Y}}
\newcommand{\bh}{\widehat{\beta}}
\newcommand{\ah}{\widehat{\alpha}}
\renewcommand{\arraystretch}{1.25}

% Небольшой запас на растяжку строк: убирает единичные overfull-строки
% в русском тексте, где переносов меньше, чем в английском.
\emergencystretch=1.5em

%
%  Компилировать XeLaTeX'ом (нужен fontspec), два прохода.
%  Драйвер обязан определить \newif\ifsolutions ДО этого \input.
% ===================================================================

\usepackage{fontspec}
\setmainfont{Times New Roman}
\setsansfont{Arial}
\setmonofont{Consolas}

\usepackage[russian]{babel}
\usepackage{amsmath,amssymb}
\usepackage{booktabs,array,multirow}
\usepackage{enumitem}
% Регистрируем русский счётчик, иначе метка label=\asbuk*) не работает.
\makeatletter
\AddEnumerateCounter{\asbuk}{\@asbuk}{щ}
\makeatother
\usepackage{xcolor}
\usepackage{framed}
\usepackage{comment}
\usepackage[a4paper,margin=25mm,headheight=15pt]{geometry}
\usepackage{fancyhdr}
\usepackage[hidelinks]{hyperref}
\usepackage{titlesec}
\usepackage{tikz}
\usetikzlibrary{calc}

% ------------------------- цвета -----------------------------------
\definecolor{accent}{RGB}{0,80,130}
\definecolor{solbar}{RGB}{0,120,90}
\definecolor{notebar}{RGB}{170,120,0}
\definecolor{xlbar}{RGB}{110,60,140}

% --------------- название семинара для колонтитула ------------------
\newcommand{\seminarname}{Семинар}
\newcommand{\setseminar}[1]{\renewcommand{\seminarname}{#1}}

% ------------------------ колонтитулы -------------------------------
\pagestyle{fancy}
\fancyhf{}
\fancyhead[L]{\small\itshape Эконометрика \textperiodcentered\ \seminarname}
\fancyhead[R]{\small\itshape \ifsolutions Решения\else Условия\fi}
\fancyfoot[C]{\small\thepage}
\renewcommand{\headrulewidth}{0.4pt}

% ------------------------ заголовки ---------------------------------
\titleformat{\section}{\large\bfseries\color{accent}}{\thesection.}{0.6em}{}
\titleformat{\subsection}{\bfseries}{\thesubsection.}{0.6em}{}
\titlespacing*{\section}{0pt}{1.4\baselineskip}{0.6\baselineskip}

% --------------------- шапка документа ------------------------------
% \seminartitle{Семинар 2}{подзаголовок}
\newcommand{\seminartitle}[2]{%
  \setseminar{#1}%
  \begin{center}
    {\LARGE\bfseries\color{accent} #1}\\[2pt]
    {\large #2}\\[4pt]
    {\itshape \ifsolutions Условия и решения \else Условия задач \fi}
  \end{center}
  \vspace{0.8\baselineskip}}

% --------------------- блок «Задача N» ------------------------------
\newcommand{\problem}[1]{%
  \bigskip\par\noindent
  {\bfseries\color{accent}\large Задача #1}\par\nopagebreak\smallskip}

% ------------------- вертикальная линейка слева ----------------------
\newenvironment{leftbarred}[1][solbar]{%
  \def\FrameCommand{{\color{#1}\vrule width 2.5pt}\hspace{10pt}}%
  \MakeFramed{\advance\hsize-\width \FrameRestore}}%
 {\endMakeFramed}

% ------------------------ «Замечание» --------------------------------
\newenvironment{remark}{%
  \par\smallskip\begin{leftbarred}[notebar]%
  \noindent{\bfseries\color{notebar}Замечание.}\ \itshape}%
 {\end{leftbarred}\par\smallskip}

% ------------------------- «В Excel» ---------------------------------
% Видно и в условиях, и в решениях: это часть задания, а не ответ.
\newenvironment{excel}{%
  \par\smallskip\begin{leftbarred}[xlbar]%
  \noindent{\bfseries\color{xlbar}В Excel.}\ }%
 {\end{leftbarred}\par\smallskip}

% Имя пункта меню / кнопки в интерфейсе Excel
\newcommand{\ui}[1]{\textsf{\small #1}}
% Формула Excel
\newcommand{\xl}[1]{\texttt{\small #1}}

% ------------------------- «Решение» ---------------------------------
% Если \solutionsfalse — всё содержимое окружения просто выбрасывается.
\ifsolutions
  \newenvironment{solution}{%
    \par\smallskip\begin{leftbarred}[solbar]%
    \noindent{\bfseries\color{solbar}Решение.}\ }%
   {\end{leftbarred}\par\smallskip}
\else
  \excludecomment{solution}
\fi

% ---------------------- сокращения для формул ------------------------
\newcommand{\TSS}{\mathrm{TSS}}
\newcommand{\ESS}{\mathrm{ESS}}
\newcommand{\RSS}{\mathrm{RSS}}
% Матожидание и дисперсия — ажурным шрифтом (blackboard bold, amssymb).
\newcommand{\E}{\mathbb{E}}   % математическое ожидание
\newcommand{\D}{\mathbb{D}}   % дисперсия
\newcommand{\V}{\D}           % то же; старое имя, им набраны тексты семинаров
\newcommand{\Cov}{\mathrm{Cov}}
\newcommand{\yh}{\widehat{Y}}
\newcommand{\bh}{\widehat{\beta}}
\newcommand{\ah}{\widehat{\alpha}}
\renewcommand{\arraystretch}{1.25}

% Небольшой запас на растяжку строк: убирает единичные overfull-строки
% в русском тексте, где переносов меньше, чем в английском.
\emergencystretch=1.5em

%
%  Компилировать XeLaTeX'ом (нужен fontspec), два прохода.
%  Драйвер обязан определить \newif\ifsolutions ДО этого \input.
% ===================================================================

\usepackage{fontspec}
\setmainfont{Times New Roman}
\setsansfont{Arial}
\setmonofont{Consolas}

\usepackage[russian]{babel}
\usepackage{amsmath,amssymb}
\usepackage{booktabs,array,multirow}
\usepackage{enumitem}
% Регистрируем русский счётчик, иначе метка label=\asbuk*) не работает.
\makeatletter
\AddEnumerateCounter{\asbuk}{\@asbuk}{щ}
\makeatother
\usepackage{xcolor}
\usepackage{framed}
\usepackage{comment}
\usepackage[a4paper,margin=25mm,headheight=15pt]{geometry}
\usepackage{fancyhdr}
\usepackage[hidelinks]{hyperref}
\usepackage{titlesec}
\usepackage{tikz}
\usetikzlibrary{calc}

% ------------------------- цвета -----------------------------------
\definecolor{accent}{RGB}{0,80,130}
\definecolor{solbar}{RGB}{0,120,90}
\definecolor{notebar}{RGB}{170,120,0}
\definecolor{xlbar}{RGB}{110,60,140}

% --------------- название семинара для колонтитула ------------------
\newcommand{\seminarname}{Семинар}
\newcommand{\setseminar}[1]{\renewcommand{\seminarname}{#1}}

% ------------------------ колонтитулы -------------------------------
\pagestyle{fancy}
\fancyhf{}
\fancyhead[L]{\small\itshape Эконометрика \textperiodcentered\ \seminarname}
\fancyhead[R]{\small\itshape \ifsolutions Решения\else Условия\fi}
\fancyfoot[C]{\small\thepage}
\renewcommand{\headrulewidth}{0.4pt}

% ------------------------ заголовки ---------------------------------
\titleformat{\section}{\large\bfseries\color{accent}}{\thesection.}{0.6em}{}
\titleformat{\subsection}{\bfseries}{\thesubsection.}{0.6em}{}
\titlespacing*{\section}{0pt}{1.4\baselineskip}{0.6\baselineskip}

% --------------------- шапка документа ------------------------------
% \seminartitle{Семинар 2}{подзаголовок}
\newcommand{\seminartitle}[2]{%
  \setseminar{#1}%
  \begin{center}
    {\LARGE\bfseries\color{accent} #1}\\[2pt]
    {\large #2}\\[4pt]
    {\itshape \ifsolutions Условия и решения \else Условия задач \fi}
  \end{center}
  \vspace{0.8\baselineskip}}

% --------------------- блок «Задача N» ------------------------------
\newcommand{\problem}[1]{%
  \bigskip\par\noindent
  {\bfseries\color{accent}\large Задача #1}\par\nopagebreak\smallskip}

% ------------------- вертикальная линейка слева ----------------------
\newenvironment{leftbarred}[1][solbar]{%
  \def\FrameCommand{{\color{#1}\vrule width 2.5pt}\hspace{10pt}}%
  \MakeFramed{\advance\hsize-\width \FrameRestore}}%
 {\endMakeFramed}

% ------------------------ «Замечание» --------------------------------
\newenvironment{remark}{%
  \par\smallskip\begin{leftbarred}[notebar]%
  \noindent{\bfseries\color{notebar}Замечание.}\ \itshape}%
 {\end{leftbarred}\par\smallskip}

% ------------------------- «В Excel» ---------------------------------
% Видно и в условиях, и в решениях: это часть задания, а не ответ.
\newenvironment{excel}{%
  \par\smallskip\begin{leftbarred}[xlbar]%
  \noindent{\bfseries\color{xlbar}В Excel.}\ }%
 {\end{leftbarred}\par\smallskip}

% Имя пункта меню / кнопки в интерфейсе Excel
\newcommand{\ui}[1]{\textsf{\small #1}}
% Формула Excel
\newcommand{\xl}[1]{\texttt{\small #1}}

% ------------------------- «Решение» ---------------------------------
% Если \solutionsfalse — всё содержимое окружения просто выбрасывается.
\ifsolutions
  \newenvironment{solution}{%
    \par\smallskip\begin{leftbarred}[solbar]%
    \noindent{\bfseries\color{solbar}Решение.}\ }%
   {\end{leftbarred}\par\smallskip}
\else
  \excludecomment{solution}
\fi

% ---------------------- сокращения для формул ------------------------
\newcommand{\TSS}{\mathrm{TSS}}
\newcommand{\ESS}{\mathrm{ESS}}
\newcommand{\RSS}{\mathrm{RSS}}
% Матожидание и дисперсия — ажурным шрифтом (blackboard bold, amssymb).
\newcommand{\E}{\mathbb{E}}   % математическое ожидание
\newcommand{\D}{\mathbb{D}}   % дисперсия
\newcommand{\V}{\D}           % то же; старое имя, им набраны тексты семинаров
\newcommand{\Cov}{\mathrm{Cov}}
\newcommand{\yh}{\widehat{Y}}
\newcommand{\bh}{\widehat{\beta}}
\newcommand{\ah}{\widehat{\alpha}}
\renewcommand{\arraystretch}{1.25}

% Небольшой запас на растяжку строк: убирает единичные overfull-строки
% в русском тексте, где переносов меньше, чем в английском.
\emergencystretch=1.5em
